\section{证明与推导（Proofs）}\label{sec:proofs}

本节每个结果都给出\emph{纸面推导}与\emph{Lean 形式证明}两种写法。

\subsection{重写与判定过程}
展开 $(x+1)^2 = x^2 + 2x + 1$（$x\in\lgZ$）：在纸上是分配律与合并同类项；
在 Lean 中 \texttt{ring} 一条 tactic 即可完成，因为该命题属于可判定的交换环等式：

\begin{lstlisting}
example : forall x : Int, (x + 1)^2 = x^2 + 2*x + 1 := by
  intro x
  ring
\end{lstlisting}

\texttt{omega} 处理 Presburger 算术，\texttt{norm_num} 处理闭数值：

\begin{lstlisting}
example (n : Nat) : n < n + 1 := by omega
example : (2 : Nat) + 2 = 4 := by norm_num
\end{lstlisting}

\subsection{归纳法}
证明 $n\cdot 0 = 0$。

\noindent\emph{纸面推导：} 对 $n$ 归纳。基例 $0\cdot0=0$ 由定义成立；
步进情形 $(n+1)\cdot0 = n\cdot0 + 0$（\texttt{Nat.succ\_mul}），
由归纳假设 $n\cdot0=0$ 得 $0+0=0$。

\begin{lstlisting}
theorem mul_zero_demo (n : Nat) : n * 0 = 0 := by
  induction n with
  | zero => rfl
  | succ n ih =>
      -- rw 重写后自动尝试 rfl，改完即闭合
      rw [Nat.succ_mul, ih]
\end{lstlisting}

列表上的同构命题 $|xs \mathbin{++} ys| = |xs|+|ys|$ 见
\texttt{Proofs/Induction.lean}。

\subsection{命题逻辑与反证法}
$\neg\neg P \leftrightarrow P$ 在经典逻辑下成立（Mathlib 默认经典）：

\begin{lstlisting}
example (P : Prop) : ~ ~ P <-> P := by
  constructor
  · intro h; by_contra hn; exact h hn
  · intro h; intro hn; exact hn h
\end{lstlisting}

\subsection{搜索：不会证就先搜}
\texttt{exact?}、\texttt{rw?}、\texttt{apply?} 在 Mathlib 里检索现成引理；
\texttt{decide} 对有限/可判定闭项直接求值。形式化工作中应优先使用搜索，
再考虑手写长证明。
